% ---- pg-500 （书页 488） ----

\begin{mvtable}
\phantomsection\label{tab:11.3}\mvtabcap{表 11.3}
\begin{tabular}{cc||ccc}
\toprule
运动 & 样本容量 & MMPI 指标 & 第一判别量 & 第二判别量 \\
\midrule
足球 & 158 & QE & 0.055 & $-0.098$ \\
篮球 & 42 & L & $-0.194$ & 0.046 \\
棒球 & 79 & F & $-0.047$ & $-0.099$ \\
赛艇 & 61 & K & 0.053 & $-0.017$ \\
击剑 & 50 & Hs & 0.077 & $-0.076$ \\
高尔夫 & 28 & D & 0.049 & 0.183 \\
体操 & 26 & Hy & $-0.028$ & 0.031 \\
曲棍球 & 28 & Pd & 0.001 & $-0.069$ \\
游泳 & 51 & Mf & $-0.074$ & $-0.076$ \\
网球 & 31 & Pa & 0.189 & 0.088 \\
田径 & 52 & Pt & 0.025 & $-0.188$ \\
摔跤 & 64 & Sc & $-0.046$ & 0.088 \\
 &  & Ma & $-0.103$ & 0.053 \\
 &  & Si & 0.041 & 0.016 \\
\bottomrule
\end{tabular}
\end{mvtable}

{\zihao{6} 资料来源：W. Morgan 和 R. W. Johnson.}

用前两个判别得分得出的组均值图示于\hyperref[fig:11.13]{图 11.13}. 这里, 虽然对均值相等的检验在 5\% 的水平下显著, 但基于 MMPI 得分的分离却不好(这是由于大样本容量).

\begin{mvfigure}
\mvinc{pg-500}{0.185,0.490,0.797,0.776}
\phantomsection\label{fig:11.13}\mvcap{图 11.13\quad 每项运动的判别均值 $\bar{y}'=[\bar{y}_1,\bar{y}_2]$}
\end{mvfigure}

然而判别系数却暗示: 第一判别量与 L 和 Ma 两个指标关系最密切, 而第二判别量则与 D 和 Pt 两个指标关系最密切. 下面给出这些研究人员所提供的解释.

第一判别量占总方差的 34.4\%, 它与 Mf 指标高度相关($r=-0.78$). 第二判别量占总方差的 18.3\%, 与 Sc, F 和 D 指标高度相关($r=0.66,0.54$ 和 $0.50$). 这些研究者认为, 第一判别量最好地代表了兴趣, 而第二判别量则反映了心理调整.
